\documentclass[level=2]{kidscode}
% Compact code boxes for this two-page sheet.
\lstnewenvironment{cs}{\lstset{style=kcpython,numbers=none,basicstyle=\ttfamily\footnotesize\setstretch{1.0}}}{}
\BeforeBeginEnvironment{cs}{\par\begin{tcolorbox}[kccode,unbreakable,left=4pt,right=2pt,top=0pt,bottom=0pt,before skip=2pt,after skip=3pt]}
\AfterEndEnvironment{cs}{\end{tcolorbox}}
\newcommand{\cshead}[1]{\par\smallskip\noindent{\large\bfseries\color{kcaccent}#1}\par\nopagebreak}
\setlength{\columnsep}{0.6cm}
% Two columns in the standard edition; one column in the dyslexia edition
% (its larger code font would otherwise wrap almost every line).
\newif\ifcstwocol \cstwocoltrue
\ifdefined\kcvariant \cstwocolfalse \fi
\setlength{\parskip}{0.2em}

\begin{document}
\kctitle{Python Cheat Sheet}{Level 2: keep this beside your keyboard}
\vspace*{-1.2cm}
\small\raggedright
\ifcstwocol\begin{multicols}{2}\raggedcolumns\fi

\cshead{Types of values}
\begin{tabular}{@{}ll@{}}
\texttt{7}, \texttt{-3} & \textbf{int}: whole number \\
\texttt{3.5}, \texttt{0.1} & \textbf{float}: decimal number \\
\texttt{"7"}, \texttt{"Nova"} & \textbf{str}: text (string) \\
\texttt{True}, \texttt{False} & \textbf{bool}: yes or no \\
\end{tabular}
\begin{cs}
n = int("7") * 2    # text -> int
label = str(42)     # int -> text
print(type(3.5))    # float
\end{cs}

\cshead{f-strings: build messages}
\begin{cs}
hero = "Nova"
score = 12
print(f"{hero} has {score} points")
print(f"{score / 5:.1f}")  # 2.4
\end{cs}

\cshead{and, or, not}
\begin{cs}
if score > 10 and lives > 0:
    print("Level up!")
if key == "q" or key == "Escape":
    running = False
if not game_over:
    move()
\end{cs}

\cshead{Functions: def and return}
\begin{cs}
def area(width, height):
    """Area of a rectangle."""
    return width * height

big = area(3, 4) + 1   # 13
\end{cs}
\texttt{return} hands a value back. \texttt{print} only shows it.
Variables made inside a function stay inside (scope).

\cshead{Lists}
\begin{cs}
colors = ["orange", "sky", "green"]
first = colors[0]    # index from 0
last = colors[-1]    # the last one
colors.append("blue")  # add at end
print(len(colors))   # 4
for color in colors:
    print(color)
if "green" in colors:
    print("green is there")
pick = random.choice(colors)
\end{cs}

\cshead{while loops}
\begin{cs}
lives = 3
while lives > 0:
    lives = lives - 1   # must change!

while True:          # ask again
    answer = input("How many? ")
    if answer.isdigit():  # a number?
        break        # leave the loop
    print("Try again")
\end{cs}
Stuck in a loop that never ends? Press the red \textbf{Stop} button.

\cshead{Nested loops (grids, pixel art)}
\begin{cs}
for row in range(3):
    for col in range(4):
        print(row, col)
\end{cs}

\cshead{Dictionaries: look up by name}
\begin{cs}
player = {"hero": "Nova", "score": 0}
player["score"] = player["score"] + 5
player["lives"] = 3         # new key
for key, value in player.items():
    print(key, value)
if "lives" in player:
    print(player["lives"])
speed = player.get("speed", 1)
\end{cs}

\cshead{Strings}
\begin{cs}
line = "  12 apples \n"
clean = line.strip()   # "12 apples"
parts = clean.split()  # 2 words
print("Hi".upper(), "Hi".lower())
\end{cs}

\cshead{Files: save and load}
\begin{cs}
with open("scores.txt", "w") as file:
    file.write("15\n")
with open("scores.txt", "a") as file:
    file.write("9\n")

try:
    with open("scores.txt") as file:
        scores = file.read().split()
except FileNotFoundError:
    scores = []      # no file yet
\end{cs}
Careful: \texttt{"w"} empties the file first!

\cshead{Testing with pytest}
\begin{cs}
def test_area_normal():
    assert area(3, 4) == 12

def test_area_edge():     # edge case
    assert area(0, 4) == 0
\end{cs}
Test names start with \texttt{test\_}. To run them, open
\texttt{my\_tests.py} (or \texttt{check\_my\_work.py}) and press Run. No pytest? \textbf{Tools > Manage packages}.

\cshead{Turtle events and games}
\begin{cs}
screen = turtle.Screen()
screen.tracer(0)       # no flicker

def go_left():
    player.setx(player.xcor() - 20)

screen.onkey(go_left, "Left")  # no ()
screen.listen()        # hear keys

def game_loop():
    if player.distance(star) < 20:
        print("caught!")   # touching
    screen.update()
    screen.ontimer(game_loop, 50)

game_loop()
turtle.done()
\end{cs}
Key names: \texttt{"Left"}, \texttt{"Right"}, \texttt{"Up"}, \texttt{"Down"},
\texttt{"space"}, \texttt{"a"}. Click the window first!

\cshead{Error message decoder}
Read the \textbf{last line} first. The line above it shows where.\\[2pt]
\begin{tabular}{@{}>{\ttfamily\scriptsize}p{3.05cm}>{\footnotesize\raggedright\arraybackslash}p{4.8cm}@{}}
SyntaxError & typo: missing \texttt{:} \texttt{)} or \texttt{"} \\
IndentationError & spaces at the start of a line \\
NameError & a name is misspelled or not made yet \\
TypeError & mixing types: \texttt{"5" + 5} \\
IndexError & no item at that index (too big) \\
KeyError & that key is not in the dictionary \\
FileNotFoundError & no file with that name in this folder \\
ValueError & \texttt{int("hello")} cannot work \\
\end{tabular}

\cshead{Debugging strategies}
\textbf{1. Read the traceback} (last line first). \textbf{2. Print-debugging:}
\texttt{print("x is", x)} to see values. \textbf{3. Rubber duck:} explain each line out loud.

\ifcstwocol\end{multicols}\fi
\end{document}
